\section{Domain Transfer: Cosplay vs. Character Fine-tuning}

Cosplay photographs are considerably more expensive to curate than character
artwork: verifying that a photo actually depicts the correct character means
checking costume accuracy, filtering out low-quality or ambiguous shots, and
often distinguishing between costume variants of the same character. Character
artwork, by contrast, can largely be collected with a single search query per
character and a quick visual pass to discard obvious mismatches. This raises
a practical question: how much recognition performance is sacrificed by
fine-tuning on the cheap, mostly-unsorted character artwork instead of the
carefully curated cosplay set?

To answer this, the 10-class subset from Table~\ref{tab:siglip2_baseline_errors}
was used to run two separate fine-tuning experiments: one trained exclusively
on \textit{cosplay photographs} (1{,}351 images, curated) and one trained
exclusively on \textit{character artwork} (1{,}626 images, mostly unsorted).
Both were evaluated on the same shared validation/test split used in the
30-class experiment, restricted to these 10 classes. That evaluation split
is itself drawn from the same crawl as the cosplay training images (matching
filename conventions and source IDs), while the character artwork comes from
a visually distinct, separately sourced crawl --- worth keeping in mind when
interpreting the results below.

\subsection{Baseline}

\begin{table}[H]
\centering
\begin{tabular}{lccl}
\toprule
\textbf{Class} & \textbf{Top-1} & \textbf{Top-5} & \textbf{Most confused as} \\
\midrule
Frieren (\textit{Frieren})            & 0.0\%  & 14.6\% & Yor Forger, Kafka (19/41 each) \\
Yor Forger (\textit{Spy x Family})    & 41.9\% & 100\%  & Kafka (17/43), Ai Hoshino (8/43) \\
Ai Hoshino (\textit{Oshi no Ko})      & 78.8\% & 100\%  & Furina (4/33), Sailor Moon (3/33) \\
Kafka (\textit{Honkai Star Rail})     & 87.1\% & 100\%  & Yor Forger (3/31) \\
\bottomrule
\end{tabular}
\caption{Classes with the lowest zero-shot top-1 accuracy on the 10-class
subset (SigLIP2, validation split). All remaining classes reached
$\geq$92\%.}
\label{tab:domain_baseline_errors}
\end{table}

The pretrained model fails on \textit{Frieren} even more severely here than
in the 30-class evaluation (0\%, never in the top-5 for 85\% of images),
confirming this is a genuine gap in the pretraining data rather than a
class-count artifact.

\subsection{Fine-tuning Setup}

Both runs used the same hyperparameters as the 30-class experiment (learning
rate $10^{-6}$, batch size 32, AMP precision, single-node 2$\times$RTX 4090)
but were trained with plain contrastive loss rather than the SigLIP loss and
without a validation split during training, since both proved unstable at
this dataset size. The cosplay run converged within 10 epochs (plateauing
from epoch 5), while the character run needed 20 epochs before a clear peak
emerged, at epoch 8, followed by a confirmed decline through epoch 20.

\begin{table}[h]
\centering
\begin{tabular}{lcccc}
\toprule
\textbf{Checkpoint} & \textbf{Val Top-1} & \textbf{Val Top-5} & \textbf{Test Top-1} & \textbf{Test Top-5} \\
\midrule
Pretrained baseline         & 71.72\% & 87.93\%  & 75.80\% & 95.02\% \\
Cosplay-trained (epoch 5)   & \textbf{97.59\%} & \textbf{100\%} & \textbf{97.51\%} & \textbf{100\%} \\
Character-trained (epoch 8) & 94.14\% & 98.97\%  & 93.95\% & 100\%  \\
\bottomrule
\end{tabular}
\caption{Validation and test accuracy before and after fine-tuning on each
domain (10-class split, 290 validation / 281 test images).}
\label{tab:domain_finetune}
\end{table}

\subsection{Results}

Both training domains produce large gains over the pretrained baseline.
Fine-tuning on cosplay photos lifts test top-1 accuracy by
\textbf{+21.7 points} (75.8\%$\to$97.5\%) and reaches 100\% top-5 on both
splits. Fine-tuning on character artwork alone lifts test top-1 accuracy by
\textbf{+18.2 points} (75.8\%$\to$94.0\%) --- despite the character training
set being 20\% larger (1{,}626 vs.\ 1{,}351 images), it falls 3.5 points
short of the cosplay-trained checkpoint on both val and test. The most
likely explanation is the domain match noted above: the validation/test
images come from the same crawl as the cosplay training set, so the
cosplay-trained model has a built-in advantage that the character-trained
model has to overcome by generalizing across a genuine visual domain gap
(illustration vs.\ photograph).

\begin{table}[H]
\centering
\begin{tabular}{lccl}
\toprule
\textbf{Class} & \textbf{Cosplay Top-1} & \textbf{Character Top-1} & \textbf{Character confused as} \\
\midrule
Frieren (\textit{Frieren})         & 87.8\%  & 73.2\%  & Yor Forger (11/41) \\
Kafka (\textit{Honkai Star Rail})  & 100\%   & 87.1\%  & Yor Forger (4/31) \\
Yor Forger (\textit{Spy x Family}) & 97.7\%  & 100\%   & --- \\
Ai Hoshino (\textit{Oshi no Ko})   & 100\%   & 100\%   & --- \\
\bottomrule
\end{tabular}
\caption{The four weakest zero-shot classes from
Table~\ref{tab:domain_baseline_errors}, compared across the two fine-tuned
checkpoints (validation split).}
\label{tab:domain_finetune_errors}
\end{table}

The gap between the two checkpoints is concentrated in exactly the classes
that were hardest zero-shot. \textit{Frieren} improves substantially under
either domain but remains 14.6 points weaker under character training
(73.2\% vs.\ 87.8\%). More strikingly, \textit{Kafka} is essentially unmoved
by character fine-tuning (87.1\%$\to$87.1\%, still confused with Yor Forger
on the same proportion of images) while cosplay fine-tuning resolves it
completely (100\%). Conversely, character training has a slight edge on
\textit{Yor Forger} (100\% vs.\ 97.7\%), so the gap is not uniform in one
direction, but cosplay training dominates on net.

\subsection{Is Curation Worth It?}

Comparing the two checkpoints directly: cosplay-trained fine-tuning reaches
97.5\% test top-1, character-trained fine-tuning reaches 94.0\% --- a gap of
3.5 points, or \textbf{83.6\%} of the achievable improvement over the
pretrained baseline (18.15 of 21.71 points) captured by character artwork
alone. Given how much cheaper character artwork is to collect --- a single
search query per character plus a quick visual pass, versus the manual
costume and quality verification cosplay photos require --- this is a
favorable trade-off for most use cases. The remaining gap is not spread
evenly across classes either: it is concentrated in a couple of already-hard
characters (\textit{Frieren}, \textit{Kafka}) rather than a uniform quality
penalty, so a hybrid strategy is plausible too: fine-tune primarily on cheap
character artwork, and only spend curation effort on cosplay photos for the
specific classes that end up weak.

In short, training on mostly-unsorted character images instead of curated
cosplay photos is a feasible substitute for this task: it recovers the large
majority of the fine-tuning gain, at a fraction of the collection cost, and
its shortfall is concentrated enough to be patched selectively rather than
requiring wholesale cosplay curation.
